\documentclass[10pt,twocolumn]{article}

\usepackage[margin=0.72in]{geometry}
\usepackage{amsmath,amssymb}
\usepackage{booktabs}
\usepackage{graphicx}
\usepackage{xcolor}
\usepackage{microtype}
\usepackage{enumitem}
\usepackage{url}
\usepackage[hidelinks]{hyperref}

\setlength{\columnsep}{0.22in}
\setlist{nosep,leftmargin=*}

\newcommand{\method}{RotQuant}
\newcommand{\todo}[1]{\textcolor{red}{[TODO: #1]}}
\newcommand{\bpw}{\mathrm{bpw}}
\newcommand{\bpv}{\mathrm{bpv}}
% Generated by scripts/audit_publication.py; do not edit by hand.
\newcommand{\QwenSourcePPL}{13.9001}
\newcommand{\QwenJointMeanPPL}{14.5548}
\newcommand{\QwenJointMeanRelativePct}{4.71}
\newcommand{\QwenJointWorstRelativePct}{5.78}
\newcommand{\EstimatedWeightReductionPct}{56.78}
\newcommand{\ActualTensorReductionPct}{59.34}
\newcommand{\ActualSnapshotReductionPct}{59.25}
\newcommand{\SourceTensorGB}{9.3197}
\newcommand{\SourceTensorExclMtpGB}{9.0785}
\newcommand{\LikeForLikeTensorReductionPct}{58.26}
\newcommand{\RotQuantTensorGB}{3.7898}
\newcommand{\CacheShortReductionPct}{24.0}
\newcommand{\CacheLongReductionPct}{10.6}
\newcommand{\MatchedCacheMeanRatio}{0.802}
\newcommand{\MatchedCacheWorstRatio}{0.835}
\newcommand{\NativeGGUFGB}{3.1463}
\newcommand{\SourceRevision}{\texttt{3764fa359b90}}
\newcommand{\RotQuantRevision}{\texttt{986ed9839cd3}}


\title{\method: Whole-System Rotation-Aware Compression of\\
Language-Model Weights and Key--Value Caches}

\author{
  Daniel Halwell\\
  Independent Researcher\\
  \texttt{\todo{email}}
  \and
  \todo{additional authors and affiliations}
}

\date{}

\begin{document}

\maketitle

\begin{abstract}
Low-bit post-training quantization is commonly evaluated on either model
weights or the key--value (KV) cache in isolation, even though both contribute
to the memory and bandwidth cost of autoregressive inference. We present
\method, a reproducible framework for jointly designing rotation-aware weight
and KV-cache compression and carrying the selected representation into a native
runtime. \method applies structured randomized Hadamard rotations, Gaussian
scalar codebooks, exact packed-bit accounting, and optional validation-gated
block reconstruction. For the cache, it searches key and value precision
independently per layer on held-out calibration calls, then freezes the selected
map for deployment and cross-context evaluation.

In a development study on Qwen3.5-4B, a uniform 4-bit weight profile obtains
mean WikiText-2 perplexity \QwenJointMeanPPL{} across three seeds, a
\QwenJointMeanRelativePct\% relative increase over matched high-precision
references. The exported checkpoint's tensor files are
\LikeForLikeTensorReductionPct\% smaller than the pinned source checkpoint's
loaded tensors, counting the same tensor set on both sides. All key--value
cache quality results in earlier drafts were measured with a cache simulator
that shared linear-attention state between its two decode passes; they are
withdrawn pending re-measurement (Section~\ref{sec:validity}), and no cache
claim is made here. Negative
results show that learned rotations, block recovery, and LoRA-based
quantization-aware recovery do not improve every model or operating point.
We additionally implement a self-describing GGUF extension and llama.cpp/Metal
prototype that preserves the packed representation rather than reconstructing
and requantizing it. The present results are development-scale: broader tasks,
models, hardware, and native end-to-end benchmarks remain required before a
state-of-the-art or production claim.
\end{abstract}

\section{Introduction}

Autoregressive model inference has two persistent memory consumers. Model
weights are read repeatedly during prefill and decoding, while the KV cache
grows linearly with sequence length and active batch size. Quantizing one while
leaving the other at half precision can simply move the memory and bandwidth
bottleneck. Moreover, the best isolated choice is not necessarily the best
whole-system choice: a slightly more accurate weight model can be more or less
sensitive to subsequent cache compression, and a cache recipe selected at one
context length may not transfer to another.

Orthogonal rotations are a promising common primitive. Randomized Hadamard
transforms make weights and activations more incoherent and suppress coordinate
outliers without changing the represented linear function
\cite{che2023quip,tseng2024quipsharp,ashkboos2024quarot}. Learned rotations can
improve over unlucky random bases under some quantization settings
\cite{liu2024spinquant}. Separately, Gaussian-aware scalar quantizers and
near-optimal vector quantizers exploit the distribution induced by rotation
\cite{dettmers2023qlora,zandieh2025turboquant}. These ideas motivate a unified
question: can one use a common rotation-and-codebook vocabulary to compress
both weights and cache state, select a joint operating point under held-out
quality gates, and deploy exactly that representation?

This paper makes four contributions:

\begin{enumerate}
  \item We define a rotation-consistent weight format combining blockwise
  randomized Hadamard transforms, MSE-searched group scales, Gaussian scalar
  codebooks, and exact storage accounting.
  \item We introduce a held-out mixed-precision cache allocator that treats key
  and value sensitivity independently at each layer, measures global teacher
  divergence per exact byte, and freezes the selected deployment map. Its
  recorded evaluation on Qwen3.5-4B is withdrawn (Section~\ref{sec:validity});
  the evaluator now fails closed unless a near-lossless 8-bit cache reproduces
  the full-precision cache on the same calls.
  \item We report a joint development study on OPT and Qwen3.5-4B, including
  multiple seeds, matched controls, cross-context transfer, negative recovery
  results, and two documented protocol corrections, the second of which
  withdraws every previously reported cache result.
  \item We implement a versioned packed artifact and native llama.cpp/Metal
  path for weights and the selected cache profile, establishing the system
  boundary needed for honest memory and performance measurement.
\end{enumerate}

The emphasis is methodological rather than a state-of-the-art claim. Our most
important finding is that compression decisions must be evaluated as deployed
systems: local reconstruction loss, nominal bit width, and even a successful
load are insufficient evidence of end-to-end quality or efficiency.

\section{Background and Related Work}

\subsection{Weight-only quantization}

GPTQ uses approximate second-order information to propagate rounding error
across columns of a weight matrix \cite{frantar2022gptq}. AWQ instead uses
activation statistics to identify and protect salient channels through an
equivalent scaling transformation \cite{lin2023awq}. AQLM formulates extreme
compression using multiple learned additive codebooks and block-level
reconstruction \cite{egiazarian2024aqlm}. These approaches demonstrate that
weight-only quantization can retain strong language-model quality, but they
occupy different points in the preprocessing, calibration, coding, and runtime
design space.

\method uses scalar rather than vector codes in the evaluated 4-bit profile.
This is deliberately simple and runtime-oriented. At 2--3 bits, vector and
trellis methods such as QuIP\#, AQLM, and QTIP provide a stronger comparison
class \cite{tseng2024quipsharp,egiazarian2024aqlm,tseng2024qtip}; we therefore
do not infer low-bit superiority from the present experiments.

\subsection{Rotation-aware quantization}

For a row activation $x\in\mathbb{R}^{d}$, weight
$W\in\mathbb{R}^{m\times d}$, and orthogonal matrix $R$, the linear map obeys
\begin{equation}
  xW^{\top} = (xR^{\top})(WR^{\top})^{\top},
  \label{eq:rotation-invariance}
\end{equation}
because $R^{\top}R=I$. QuIP and QuIP\# use incoherence processing to make
weight and Hessian directions less aligned with coordinate axes
\cite{che2023quip,tseng2024quipsharp}. QuaRot extends function-preserving
rotations to weights, activations, and the cache \cite{ashkboos2024quarot},
while SpinQuant learns rotations using task-aware optimization
\cite{liu2024spinquant}. TurboQuant studies near-optimal distortion after random
rotation and introduces a residual one-bit randomized sketch for unbiased inner
products \cite{zandieh2025turboquant}.

\method builds on these principles but focuses on a distinct systems problem:
the selection, validation, serialization, and native execution of one joint
weight-and-cache profile. The evaluated winner does not use TurboQuant's
residual QJL sketch; it uses a Gaussian scalar codebook and groupwise scale
search after structured rotation.

\subsection{KV-cache quantization}

KIVI observes different distributional structure for keys and values and uses
asymmetric 2-bit cache quantization \cite{liu2024kivi}. KVQuant combines
per-channel keys, non-uniform datatypes, and outlier handling
\cite{hooper2024kvquant}. These results establish that keys and values should
not automatically share one quantization rule. \method extends this principle
to per-layer rate allocation and evaluates a frozen map across context lengths.

\section{Method}

Figure~\ref{fig:overview} summarizes the pipeline. A source model is first
profiled under candidate weight recipes. Viable weight profiles are then
crossed with uniform and mixed cache profiles. Selection calls, evaluation
calls, and final perplexity data are disjoint. The chosen representation is
serialized without reconstructing dense weights.

\begin{figure*}[t]
\centering
\fbox{\parbox{0.96\textwidth}{\centering
\textbf{Source model}
$\rightarrow$ rotation-consistent weight candidates
$\rightarrow$ held-out weight gate
$\rightarrow$ per-layer K/V sensitivity and exact-byte allocation
$\rightarrow$ frozen joint profile
$\rightarrow$ multi-seed and cross-context validation
$\rightarrow$ packed artifact
$\rightarrow$ native runtime\\[2pt]
Selection data, evaluation data, and final quality data remain disjoint;
failed recovery stages restore the previous deployable checkpoint.}}
\caption{The \method selection and deployment pipeline. This placeholder will
be replaced by a graphical system diagram in the next draft.}
\label{fig:overview}
\end{figure*}

\subsection{Structured rotation}

We use a randomized block Hadamard rotation
\begin{equation}
  R = \operatorname{blockdiag}\!\left(\frac{H_b}{\sqrt{b}}\right)D,
\end{equation}
where $H_b$ is a Walsh--Hadamard matrix and $D$ is a fixed random sign
diagonal. The evaluated Qwen profile uses block size $b=128$. This transform is
exactly orthogonal and costs $O(d\log b)$ additions per vector. If a dimension
is not divisible by 128, the implementation selects the largest dividing
power-of-two block.

The framework also provides an exactly orthogonal trainable butterfly. Each
Hadamard stage is replaced by pairwise transforms parameterized by angles
$\theta$ and initialized at $\pi/4$, making initialization numerically
equivalent to the randomized Hadamard transform. This parameterization enables
validation-gated learning while retaining $O(d\log b)$ complexity. Learned
rotation is an ablation, not part of the final Qwen profile.

Rotation consistency is enforced per linear operation: the activation and
input dimension of the weight use the same basis. A deliberately mismatched
control rotates weights without the corresponding activation transformation;
this exposes cross-layer drift and is never exportable.

\subsection{Gaussian scalar weight quantization}

After rotation, each row of $\widetilde W=WR^{\top}$ is split into groups of
$g$ input coordinates. Let $C_b=\{c_1,\ldots,c_{2^b}\}$ be a scalar Lloyd--Max
codebook optimized for a standard Gaussian source. For group $G$, we select a
positive scale $s_G$ from a fixed search grid around its RMS value:
\begin{equation}
  s_G^* = \arg\min_{s\in\mathcal{S}(G)}
  \sum_{w\in G}\left(w-s\,Q_{C_b}(w/s)\right)^2.
  \label{eq:scale-search}
\end{equation}
Scales are rounded to FP16 before assigning final indices, ensuring that
pack-time and runtime reconstruction agree. With $b$ code bits, FP16 scales,
and full groups of size $g$, the logical rate is
\begin{equation}
  r_W = b + \frac{16}{g}\quad\bpw.
  \label{eq:weight-rate}
\end{equation}
Thus the evaluated 4-bit, group-128 profile uses 4.125 logical bits per
quantized weight, excluding unquantized model components and small format
metadata. We report complete-model estimates separately from quantized-tensor
rates.

\subsection{Validation-gated block recovery}

Local weight MSE does not capture error propagation through a transformer.
\method can replay complete transformer blocks on recorded source calls,
initialize trainable butterflies from the deployable randomized rotation, and
optionally learn group scale multipliers. Later blocks can receive the already
quantized preceding state to expose accumulated drift. Candidate checkpoints
are selected on disjoint block calls; if no checkpoint beats step zero, the
source deployable profile is restored.

An optional LoRA recovery stage \cite{hu2021lora} optimizes low-rank residuals
against teacher logits while retaining exact adapter-byte accounting. The same
rollback rule applies. In the Qwen experiments, every LoRA arm selected step
zero, and recovery on the final winner slightly worsened perplexity; neither is
included in the released profile.

\subsection{Rotated KV-cache quantization}

Let post-RoPE queries, keys, and values for one head be $Q,K,V$. We use
independent orthogonal rotations $R_K$ and $R_V$:
\begin{align}
  \widetilde Q &= QR_K^{\top}, &
  \widetilde K &= KR_K^{\top}, &
  \widetilde V &= VR_V^{\top}.
\end{align}
The attention score is preserved before quantization because
$\widetilde Q\widetilde K^{\top}=QK^{\top}$. After multiplying attention
probabilities by the dequantized $\widetilde V$, the output is transformed by
$R_V$ to return to the original basis.

Cache rows use Gaussian scalar codes with group size 64 and FP16 group scales.
For code width $b$, the logical rate is therefore $b+16/64=b+0.25$ bits per
stored value. The deployed map chooses $b_K,b_V\in\{2,3,4\}$ independently for
each full-attention layer, with average rate 3.25~bpv including scales.

\subsection{Held-out rate allocation}

The allocator evaluates one layer-side candidate at a time on selection calls.
Its objective is global teacher KL reduction per exact byte relative to a
lower-rate state. It constructs a recipe under a target byte budget, measures
the assembled recipe jointly, and falls back to the best same-budget uniform
profile if interactions overturn the predicted gain. Candidate construction
uses no final evaluation calls. Before any candidate is scored, a uniform
8-bit Gaussian cache must reproduce the full-precision cache on the same
held-out calls (teacher KL below $10^{-2}$); otherwise the run is rejected,
because a floor that survives 8-bit codes cannot be quantization error.

The final Qwen map contains eight full-attention layers. Mean selected
key/value code widths are not constrained to be equal. A mixed-context map is
constructed using both short- and long-context selection metrics, then frozen
and replayed without recalibration on held-out calls at both lengths.

\subsection{Artifact and native runtime}

The Python checkpoint stores packed codes, FP16 scales, codebook identifiers,
rotation parameters, logical shapes, model configuration, tokenizer assets,
and deployment metadata without pickle. A native GGUF extension records the
same semantics under a versioned \texttt{rotquant.*} namespace. The Qwen
prototype stores 4-bit Gaussian weight rows and the exact 2/3/4-bit frozen
cache map. It also represents the tied input embedding and output projection
once.

The audited exported Transformers checkpoint contains 200 quantized modules,
no retained LoRA adapter, and \RotQuantTensorGB{}~GB of tensor files. The
pinned source checkpoint holds \SourceTensorGB{}~GB of tensors, of which
0.241~GB is a multi-token-prediction head (\texttt{mtp.*}) that the
Transformers model class never loads and the export therefore never stores.
Against the \SourceTensorExclMtpGB{}~GB of source tensors that are actually
loaded, the export is a measured \LikeForLikeTensorReductionPct\%
tensor-storage reduction; the whole-file figures
(\ActualTensorReductionPct\% for tensor files, \ActualSnapshotReductionPct\%
for complete snapshots) overstate the reduction because they count the dropped
head on the source side only. The native language-path GGUF is
\NativeGGUFGB{}~GB, but it omits the vision tower and is therefore not used for
the like-for-like storage comparison.

The llama.cpp fork provides a CPU reference and Metal kernels for structured
weight rotation and packed Gaussian-codebook matrix--vector products. The
cache is written in packed form and retained persistently; the current Metal
prototype dequantizes each active cache block into transient graph storage
before flash attention. This is native representation preservation, but not yet
a fully fused cache-attention kernel.

\section{Experimental Methodology}

\subsection{Models and data}

Early ablations use OPT-125M and OPT-1.3B \cite{zhang2022opt}. The principal
study uses the language path of Qwen3.5-4B. \todo{Add the exact upstream model
revision and an authoritative Qwen3.5 architecture citation.} Weight quality is
measured on WikiText-2 \cite{merity2016wikitext}; cache selection and held-out
trajectory comparisons use disjoint C4 calls \cite{raffel2020t5,dodge2021c4}.

The main Qwen protocol evaluates 64 sequences of length 256 for perplexity.
Cache experiments use separate selection and evaluation batches and generate
16 or 32 continuation tokens depending on the stage. Long-context confirmation
uses four 1,024-token selection calls and four disjoint evaluation calls.

\subsection{Metrics}

We report perplexity and relative perplexity against a matched high-precision
reference. Cache fidelity includes teacher KL divergence, cosine similarity,
top-1 token agreement, negative-log-likelihood delta, and direct K/V
reconstruction NMSE. Rates include packed codes and FP16 scale bytes. Estimated
complete-model size includes unquantized components; it is labeled as an
estimate until replaced by a verified, like-for-like artifact comparison.

Each cache KL uses the same weight model with a full-precision cache as its
teacher. Consequently, cache KL compares cache sensitivity within a weight
profile; it is not a cross-weight accuracy metric. Likewise, a lower cache KL
for quantized weights than source weights does not imply that quantized weights
are more accurate.

\subsection{Protocol safeguards}

The evaluation code requires:

\begin{itemize}
  \item matched model, tokenizer, split, sequence length, and sample count;
  \item disjoint calibration, candidate-selection, and final evaluation calls;
  \item exact byte targets and endpoint-equivalence tests;
  \item at least three seeds for promoted candidates;
  \item step-zero rollback for learned recovery;
  \item explicit labeling of dense-weight fallback as quality-only; and
  \item matched cache controls within weight profile and seed.
\end{itemize}

An early cache matrix violated the second requirement: uniform and dynamic rows
used different C4 calls. The discrepancy was revealed because nominally
identical uniform-8-bit and restricted-dynamic-8-bit endpoints disagreed. We
discarded the cross-profile ranking, corrected the evaluation offset, and added
an endpoint-equivalence regression test. Only corrected runs appear below.

\subsection{Validity notice: cache results withdrawn}
\label{sec:validity}

A second, more serious protocol failure was found after the cache study was
complete. The Transformers cache simulator cloned only tensor-valued cache
attributes. The Transformers release that the cache notebooks installed
(5.16.x) keeps each linear-attention layer's convolution and recurrent state in
dictionary attributes that are updated in place, so the ``source'' and
``packed'' decode passes shared that state and corrupted each other. The
resulting next-token KL was independent of the code width: uniform 8-bit,
4-bit, and 2-bit caches all reported KL of roughly 0.55, and 4-bit was worse
than 2-bit in two of three seeds. Replicating the protocol on the pinned
Qwen3.5-4B weights with the corrected simulator gives KL of $5\times10^{-4}$,
$6\times10^{-3}$, and $6.5\times10^{-2}$ for 8-, 4-, and 2-bit caches, whereas
the defective simulator under the same Transformers release reproduces the
recorded floor ($0.88$ at 8 bits).

Every cache KL previously reported by this draft, including the
uniform-versus-mixed tables, the cross-context transfer table, the matched
K4/V4 controls, and the selection of the frozen 3.25-bpv map, is therefore
withdrawn. Tables~\ref{tab:cache-transfer} and~\ref{tab:matched} are retained
only as a record of the withdrawn numbers. Weight-only perplexity results do
not use the simulator and stand. The simulator now clones all cache state,
fails closed on shared storage, and evaluates a uniform 8-bit cache on the same
held-out calls before any candidate; a run whose 8-bit KL exceeds $10^{-2}$ is
rejected. The cache study will be repeated under that protocol before any
cache claim is made.

\section{Results}

\subsection{Rotation is necessary but not sufficient}

Table~\ref{tab:opt} shows the directionally stable role of rotation at 3 bits.
On OPT-1.3B, randomized Hadamard rotation reduces perplexity from 313.53 to
58.89, although it remains well above the 30.33 source reference. On OPT-125M,
validation-gated scale learning and propagation-aware block inputs improve over
the basic rotation, while end-to-end distillation makes no change.

\begin{table}[t]
\centering
\small
\caption{Development perplexity on WikiText-2. Values after the first two
OPT-125M rows are single-seed ablations. The learned stages are cumulative.}
\label{tab:opt}
\resizebox{\columnwidth}{!}{%
\begin{tabular}{lrr}
\toprule
Model / profile & Bits & PPL $\downarrow$ \\
\midrule
OPT-125M, no rotation & 3.125 & 793.3127 \\
OPT-125M, randomized FWHT & 3.125 & $112.2128\pm1.3688$ \\
\quad learned butterfly & 3.125 & 107.7451 \\
\quad learned group scales & 3.125 & 92.8261 \\
\quad propagated inputs & 3.125 & 83.2201 \\
\quad LoRA-QAT recovery & 3.125+adapter & 76.0044 \\
\midrule
OPT-1.3B, source & 16 & 30.3330 \\
OPT-1.3B, no rotation & 3.125 & 313.5290 \\
OPT-1.3B, randomized FWHT & 3.125 & $58.8901\pm4.2042$ \\
\bottomrule
\end{tabular}
}
\end{table}

These results also expose a limitation of local optimization. Increasing
butterfly steps from 12 to 18 did not improve deployed perplexity, and several
blocks overfit local reconstruction data. Learning deployable scale/clipping
parameters helped more than additional rotation steps.

\subsection{Four-bit weights are the viable Qwen operating point}

Table~\ref{tab:qwen-weights} reports the seed-0 joint-matrix screen. Uniform
4-bit weights are both more accurate and nearly the same estimated size as the
dynamic 4.125-bpw allocation. Uniform 3-bit weights fail the predeclared 10\%
quality gate. In an earlier matched matrix, rank-4 and rank-8 LoRA recovery both
selected step zero; the 3-bit rank-4 arm reached 24.39\% relative perplexity
degradation for only 4.61 additional percentage points of estimated total-size
reduction over the 4-bit control.

\begin{table*}[t]
\centering
\small
\caption{Qwen3.5-4B seed-0 weight screen. Complete sizes are logical estimates
that include unquantized components. Relative perplexity is measured against
the matched 13.9001 source run.}
\label{tab:qwen-weights}
\resizebox{\textwidth}{!}{%
\begin{tabular}{lrrrr}
\toprule
Profile & Quantized bpw & PPL $\downarrow$ & Relative PPL & Estimated complete weights \\
\midrule
Source & 16 & 13.9001 & 0.00\% & 9.3200 GB \\
Uniform W4 & 4.125 & \textbf{14.7029} & \textbf{+5.78\%} & 4.0277 GB ($-56.78\%$) \\
Dynamic, 4.125-bpw target & 4.116 & 14.9976 & +7.90\% & \textbf{4.0238 GB} ($-56.83\%$) \\
\bottomrule
\end{tabular}
}
\end{table*}

Dynamic weights save only 3.93 MB (0.098\%) relative to uniform W4 while
increasing perplexity by a further 2.00\%. This illustrates why a layerwise
sensitivity search should not be retained merely because it returns a valid
mixed recipe.

\subsection{Frozen cache-map transfer (withdrawn)}

\textbf{Withdrawn (Section~\ref{sec:validity}).}
Table~\ref{tab:cache-transfer} compared three frozen mixed maps with uniform
K3/V3 at the identical 3.25-bpv storage rate and reported KL reductions of
\CacheShortReductionPct\% at 256 tokens and \CacheLongReductionPct\% at
1,024 tokens relative to uniform K3/V3. These values were produced by the
defective simulator and do not measure quantization error; the table is kept
only as a record.

\begin{table}[t]
\centering
\small
\caption{Withdrawn: frozen 3.25-bpv cache-map transfer as measured by the
defective simulator (Section~\ref{sec:validity}). Retained as a record only.}
\label{tab:cache-transfer}
\begin{tabular}{lrr}
\toprule
Frozen map & 256-token KL & 1,024-token KL \\
\midrule
Short-context map & 0.445652 & 1.430826 \\
Long-context map & 0.496695 & \textbf{1.359091} \\
Mixed-context map & \textbf{0.435155} & 1.376698 \\
Uniform K3/V3 & 0.572441 & 1.540281 \\
\bottomrule
\end{tabular}
\end{table}

The apparent reversal of key/value emphasis between the 256-token
(K2.875/V3.125) and 1,024-token (K3.375/V2.625) allocations was derived from
the same withdrawn measurements and must not be read as evidence that
key/value sensitivity changes with context length.

\subsection{Joint profile and matched controls}

Across three seeds, uniform W4 weights obtain perplexities 14.7029, 14.6091,
and 14.3524 (weight perplexity does not depend on the cache recipe). The mean
is \QwenJointMeanPPL{}, with mean relative degradation
\QwenJointMeanRelativePct\% and worst-seed degradation
\QwenJointWorstRelativePct\%. It passes the predeclared 5\% mean and 10\%
worst development gates.

\textbf{Withdrawn (Section~\ref{sec:validity}).} An initially proposed cache
gate compared the worst candidate seed with a seed-0 control and was therefore
inconclusive; missing K4/V4 controls were run at seeds 1 and 2 and replaced it
with within-seed ratios. Table~\ref{tab:matched} recorded the 3.25-bpv
candidate below the 4.25-bpv K4/V4 control in every seed, with 16.5--23.2\% KL
reduction. Those measurements came from the defective simulator and are
withdrawn together with the selection of the map itself.

\begin{table}[t]
\centering
\small
\caption{Withdrawn: matched cache comparison as measured by the defective
simulator (Section~\ref{sec:validity}). Retained as a record only.}
\label{tab:matched}
\begin{tabular}{rrrr}
\toprule
Seed & Mixed 3.25 KL & K4/V4 KL & Reduction \\
\midrule
0 & 0.435155 & 0.541733 & 19.7\% \\
1 & 0.700824 & 0.839113 & 16.5\% \\
2 & 0.514957 & 0.670372 & 23.2\% \\
\bottomrule
\end{tabular}
\end{table}

The recorded candidate/control KL ratio averaged \MatchedCacheMeanRatio{} with
a maximum of \MatchedCacheWorstRatio{}. Under the corrected simulator the
K4/V4 control itself costs only about $6\times10^{-3}$ nats on the source
model, so differences between 3.25-bpv recipes will be at the $10^{-3}$ level
and will need far more than 64 evaluated tokens per seed to resolve.

\subsection{Recovery can be actively unhelpful}

Block-and-scale recovery applied to the actual joint winner changes seed-0
perplexity from 14.702918 to 14.713017, a 0.0687\% degradation, while adding an
estimated 8.86 MB. The predeclared improvement gate rejects it. LoRA-QAT on a
different finalist and the earlier Qwen matrix also selects step zero. These
negative results support a fail-closed rule: trainable recovery must change a
held-out objective and improve the deployed metric after accounting for its
bytes.

\subsection{Preliminary native measurements}

The native weight path executes successfully on Apple Metal. Preliminary
llama.cpp measurements on an Apple M5 Max are shown in
Table~\ref{tab:native}. The quantized-cache rows use llama.cpp's generic Q4\_0
cache rather than the final Gaussian mixed profile; they are diagnostic and do
not establish a speedup. At depth 512, generic Q4 cache decoding is slower than
FP16, suggesting that dequantization overhead exceeds bandwidth savings at this
operating point.

\begin{table}[t]
\centering
\small
\caption{Preliminary native throughput in tokens/s. These data do not benchmark
the final mixed Gaussian cache kernel.}
\label{tab:native}
\begin{tabular}{llr}
\toprule
Cache & Test & Throughput \\
\midrule
FP16 & prompt, 64 tokens & $20.57\pm0.27$ \\
Q4\_0 & prompt, 64 tokens & $22.59\pm0.77$ \\
FP16 & decode 16, depth 512 & $14.72\pm2.88$ \\
Q4\_0 & decode 16, depth 512 & $12.13\pm0.45$ \\
\bottomrule
\end{tabular}
\end{table}

\section{Toward Architecture-Generic RotQuant}

The present evidence covers decoder-only language paths, but the implementation
should generalize by operator and state substrate rather than by an ever-growing
list of model class names. We propose the following abstraction:

\begin{center}
\small\ttfamily
\begin{tabular}{@{}l@{}}
PipelineAdapter \\
\quad ComponentAdapter[] \\
\quad OperatorAdapter[] \\
\quad RotationRegion[] \\
\quad StateAdapter[] \\
\quad CalibrationProfile \\
\quad EvaluationProfile \\
\quad RuntimeExportProfile
\end{tabular}
\end{center}

Support should be reported at three levels: (1) pipeline-compatible, where
unsupported components remain at high precision; (2) kernel-optimized, where
dedicated kernels cover the parameter- and latency-dominant path; and (3)
end-to-end certified, where task quality, state correctness, memory, latency,
and portability pass a modality-specific profile.

Near-term families with substantial operator reuse include dense encoders,
encoder--decoders, embedding models, MoE models, hybrid state-space models,
modular vision-language models, and Whisper-style ASR. Diffusion transformers,
streaming audio, video, robotics, and graph/spectral scientific models require
new calibration dimensions, persistent-state contracts, and kernels. This is a
research roadmap, not evidence that the current prototype supports these
families.

\section{Limitations and Threats to Validity}

\paragraph{Development-scale quality evaluation.}
The principal result uses one 4B model, a WikiText-2 subset, short continuation
trajectories, and three seeds. It lacks C4 perplexity, broad zero-shot tasks,
reasoning, code, multilingual evaluation, retrieval, and long-context
perplexity. WikiText-2 alone is not sufficient for a publication-quality model
comparison.

\paragraph{Incomplete baseline matrix.}
The draft does not yet include matched GPTQ, AWQ, QuaRot, SpinQuant, QuIP\#,
AQLM, or standard GGUF baselines on Qwen3.5-4B. The current result therefore
supports an internal method comparison, not state-of-the-art superiority.

\paragraph{Storage is measured; resident memory is not.}
The exported Transformers tensor files measure \RotQuantTensorGB{}~GB versus
\SourceTensorExclMtpGB{}~GB of loaded source tensors, a
\LikeForLikeTensorReductionPct\% reduction; the raw file-to-file figure of
\ActualTensorReductionPct\% counts the unloaded multi-token-prediction head on
the source side only. This supersedes the earlier
\EstimatedWeightReductionPct\% logical estimate for storage reporting.
Publication still requires resident and peak transient-memory measurements.
The native GGUF omits the vision tower, so its file size cannot serve as the
like-for-like multimodal comparison.

\paragraph{Cache quality is unmeasured.}
Every cache KL in earlier drafts is withdrawn (Section~\ref{sec:validity}).
The frozen map embedded in the released artifact is unvalidated until the
cache study is repeated with the endpoint-checked simulator.

\paragraph{Fallback is not deployment.}
CUDA and MPS quality experiments cache reconstructed FP16 weights for speed.
They demonstrate numerical quality but not packed runtime memory or throughput.
Only the native packed path can support deployment-performance claims.

\paragraph{Cache execution is not fully fused.}
The native cache stays packed persistently but currently dequantizes into
transient graph storage before flash attention. A fused write path and fused
dequantization-attention kernel are required to realize the expected bandwidth
advantage and to avoid transient-memory ambiguity.

\paragraph{Architecture coverage.}
Qwen3.5-4B is multimodal, but only its language path is quantized and evaluated.
Small linear-attention gates and the vision tower remain at source precision.
The paper must not characterize the present result as full multimodal support.

\section{Conclusion}

\method treats low-bit inference as a joint selection-and-deployment problem.
Structured rotations and Gaussian scalar codes make a simple 4-bit weight
profile viable, and a native format preserves its representation through
loading and execution. The cache half of the study must be repeated: its
recorded results were produced by a simulator defect that no recipe comparison
could have exposed without a near-lossless endpoint, which the protocol now
requires. Equally important, several plausible recovery and dynamic-weight
strategies fail their gates, demonstrating the value of rollback, matched
controls, and endpoint checks. The next stage is not to add more optimization
knobs; it is to re-run the cache study and the missing baseline, task,
long-context, resident-memory, and fused-runtime evaluations needed to
establish external validity.

\section*{Reproducibility Statement}

The repository records experiment configurations, random seeds, environment
metadata, exact bit accounting, negative results, protocol corrections, packed
checkpoint logic, a GGUF exporter, conformance tests, and the pinned llama.cpp
patch. The current manifest pins source revision \SourceRevision{}, RotQuant
revision \RotQuantRevision{}, and the native artifact SHA-256. Numerical claims
and local artifact bytes are checked by
\texttt{scripts/audit\_publication.py}.
The pinned GPU command matrix is generated by
\texttt{scripts/run\_publication\_suite.py}. \todo{Add public repository URL,
immutable release tag, remaining hardware identifiers, and completed GPU result
artifacts before submission.}

\bibliographystyle{unsrt}
\bibliography{references}

\appendix

\section{Frozen Qwen3.5-4B Cache Map}

The evaluated 3.25-bpv Gaussian profile uses group size 64 and the following
code widths in full-attention layers:

\begin{table}[h]
\centering
\small
\begin{tabular}{rrrrrrrr}
\toprule
Layer & 3 & 7 & 11 & 15 & 19 & 23 & 27 \\
\midrule
$K$ bits & 2 & 4 & 3 & 4 & 3 & 3 & 4 \\
$V$ bits & 2 & 4 & 3 & 2 & 3 & 2 & 3 \\
\bottomrule
\end{tabular}

\vspace{4pt}

\begin{tabular}{rr}
\toprule
Layer & 31 \\
\midrule
$K$ bits & 2 \\
$V$ bits & 4 \\
\bottomrule
\end{tabular}
\caption{Frozen mixed-context deployment map as embedded in the released
artifact. It was selected under the withdrawn measurements
(Section~\ref{sec:validity}) and is unvalidated. Adding 0.25 bpv for FP16
scales amortized over groups of 64 gives 3.25 effective bpv.}
\end{table}

\section{Publication Completion Checklist}

The following items are intentionally left open rather than represented as
completed evidence:

\begin{enumerate}
  \item Run full WikiText-2 and C4 perplexity plus a fixed zero-shot suite on
  source, RotQuant, and all matched baselines.
  \item Add at least two more model families and two model scales, including a
  dense standard architecture without custom multimodal wrappers.
  \item Evaluate long-context perplexity and retrieval at multiple cache depths.
  \item Report process resident memory, peak transient memory, prefill
  throughput, decode throughput, and energy where available. Artifact and
  tensor-file bytes are now audited by the publication manifest.
  \item Benchmark the exact Gaussian mixed cache against FP16 and standard
  cache formats after kernel fusion.
  \item Add confidence intervals and statistical tests appropriate to each
  promoted quality metric.
  \item Replace the overview placeholder with a publication-quality figure and
  add rate--quality and context-depth plots generated from raw result JSON.
  \item Archive configs, logs, checksums, source/model revisions, and the pinned
  runtime patch under one immutable release.
\end{enumerate}

\end{document}
